% !TEX root = relatorio.tex
%
% Capítulo 7 — Discussão
%
\chapter{Discussão} \label{cap:discussao}

Este capítulo analisa criticamente os resultados obtidos, confrontando-os com os objetivos propostos, identificando as decisões técnicas mais bem-sucedidas, reconhecendo as limitações encontradas, posicionando o trabalho face ao estado da arte e traçando direções para trabalho futuro.

% ─────────────────────────────────────────────────────────────────────────────
\section{Análise dos Resultados face aos Objetivos} \label{sec:analise-objetivos}

O projeto foi orientado por cinco objetivos operacionais, cuja concretização se analisa de seguida.

\textbf{OBJ-1 — Plataforma de aprendizagem adaptativa alinhada com o DigComp~2.2.} O sistema implementa todos os elementos estruturantes de uma plataforma adaptativa: inscrição com gating de pré-requisitos, entrega de tarefas por calendário, percurso de dificuldade com ramos de remediação, e escalamento para tutoria IA. O alinhamento com o DigComp~2.2 é assegurado ao nível dos tópicos, curados pelos administradores a partir de fontes documentais alinhadas com as áreas de competência da framework. A avaliação de usabilidade (§\ref{sec:testes-usabilidade}) confirmou que 66,7\% dos participantes classificaram a aplicação na categoria SUS \textit{Bom} ou \textit{Excelente}, com mediana de 85,0.

\textbf{OBJ-2 — Pipeline de geração de conteúdo por IA a partir de fontes arbitrárias.} O pipeline (§\ref{sec:rag}) funciona de ponta a ponta --- \textit{ingestion}, \textit{chunking}, geração por Mistral com \textit{schema reminder}, e deduplicação semântica --- e é robusto a falhas intermediárias: degradação graciosa por chunk e circuit breaker Resilience4j (§\ref{sec:thread-pool-cb}). O objetivo foi atingido com seis garantias de qualidade verificáveis (§\ref{sec:garantias-qualidade}).

\textbf{OBJ-3 — Remediação adaptativa personalizada por padrão de erro.} O \texttt{ErrorPatternClassifier} classifica cada resposta errada em três padrões e injeta esse diagnóstico no prompt do ramo adaptativo (§\ref{sec:consolidation}); a reutilização semântica de ramos (§\ref{sec:recuperacao}) acrescenta eficiência económica ao evitar invocações desnecessárias do modelo. O objetivo foi atingido; a limitação identificada é a profundidade máxima de \textit{struggle} (\texttt{MAX\_STRUGGLE\_DEPTH~=~1}).

\textbf{OBJ-4 — Tutoria conversacional com IA como escalamento da remediação.} A \texttt{ExplanationSession} (§\ref{sec:explicacao-ia}) implementa um ciclo completo --- geração assíncrona com contexto relacional, diálogo multi-turno e resolução com retorno ao fluxo principal --- com injeção de cabeçalho de contexto para prevenir injeção de prompt. O objetivo foi atingido; a avaliação qualitativa (§\ref{sec:usab-qualitativa}) confirmou que esta foi a funcionalidade mais valorizada.

\textbf{OBJ-5 — Segurança e qualidade do conteúdo gerado numa aplicação web multi-utilizador.} A análise de risco OWASP (§\ref{sec:riscos}) identificou dez ameaças, com os três riscos residuais de nível Médio mitigados por controlos implementados; a suite de segurança (§\ref{sec:testes-seguranca}) valida SSRF, XSS, magic link e controlo de acesso ao Swagger. O objetivo foi atingido para o perfil de ameaças identificado.

% ─────────────────────────────────────────────────────────────────────────────
\section{O que Correu Bem} \label{sec:positivo}

\textbf{Arquitetura hexagonal como base de testabilidade.} Modelar todas as dependências externas como interfaces revelou-se a decisão mais impactante do projeto: permitiu uma suite de 49 ficheiros de teste em que os unitários executam em memória pura, em microssegundos. O contrato \texttt{Either<DomainError,~T>} da Arrow tornou os casos de insucesso tão explícitos e verificáveis como os de sucesso.

\textbf{Reutilização semântica de ramos adaptativos.} O mecanismo de três estratégias (FULL\_REUSE, PARTIAL\_REUSE, FRESH\_GENERATION) por similaridade coseno pgvector é o contributo técnico mais original do sistema: servir um ramo sem qualquer chamada ao modelo (FULL\_REUSE, similaridade $> 0{,}90$) reduz os custos de API de forma estrutural à medida que a base de ramos cresce --- uma melhoria contínua sem custo marginal crescente, valiosa para uma plataforma em escala.

\textbf{Resiliência por design no pipeline de geração.} A combinação de circuit breaker, degradação graciosa por chunk, \textit{schema reminder} automático e \textit{fail-safe} na \texttt{ExplanationSession} tornou o sistema tolerante a falhas intermitentes da API externa sem expor o utilizador a erros, com visibilidade operacional imediata via alerta \texttt{AiCircuitBreakerOpen}.

\textbf{Autenticação sem senha com garantias de segurança formais.} O fluxo \textit{magic link} eliminou a complexidade de gestão de passwords e é mais seguro para utilizadores não técnicos; hashing SHA-256, consumo atómico e rotação de família de tokens cobrem as ameaças mais relevantes do modelo de autenticação (§\ref{sec:autenticacao}).

\textbf{Rate limiting com teste determinístico.} A substituição do \texttt{TimeMeter} por \texttt{ManualTimeMeter} nos testes de integração (§\ref{sec:testes-rate-int}) permitiu verificar limites e resets de janela sem dependência de tempo de parede, tornando os testes 100\% determinísticos --- uma técnica diretamente transferível a qualquer componente sensível ao tempo.

% ─────────────────────────────────────────────────────────────────────────────
\section{Limitações e Dificuldades} \label{sec:limitacoes}

\textbf{Conformidade de schema do modelo de linguagem.} O Mistral não garante que a resposta JSON respeite sempre o schema pedido. O \textit{schema reminder} mitiga o problema com uma segunda invocação, mas não elimina a possibilidade de falha persistente; sem dados de produção em escala não é possível quantificar a taxa de ocorrência real.

\textbf{Profundidade máxima de struggle.} \texttt{MAX\_STRUGGLE\_DEPTH~=~1} significa que após um único ramo falhado o sistema escala imediatamente para explicação. Numa plataforma com tópicos de nível avançado, uma profundidade maior poderia ser pedagogicamente mais adequada; a constante é configurável apenas por código, dificultando a experimentação.

\textbf{Bugs de interface identificados na avaliação.} A avaliação (§\ref{sec:usab-qualitativa}) identificou uma \textit{race condition} visual no caminho de consolidação, responsável pelas duas pontuações SUS mais baixas (P04~=~45, P05~=~35), além da barra de pesquisa oculta durante a inserção de e-mail e inconsistência de idioma nos conteúdos --- problemas que afetam a percepção de qualidade de forma desproporcionada à sua complexidade técnica de resolução.

\textbf{Ausência de métricas de produção em escala.} O sistema instrumenta 18 métricas e oito alertas (§\ref{sec:metricas}), mas a avaliação foi conduzida com 12 participantes num ambiente controlado; latências reais, distribuição de estratégias de reutilização e taxa de deduplicação efetiva sob carga real permanecem por medir.

\textbf{Dependência de fornecedor de IA único.} O sistema está acoplado ao Mistral AI. A abstracção LangChain4j isola o código de negócio da implementação concreta, mas os prompts estão calibrados ao comportamento específico do Mistral; uma migração exigiria recalibração e re-embedding de todos os vectores existentes.

\textbf{Validação pedagógica do alinhamento DigComp~2.2.} O alinhamento é garantido pela curação humana dos tópicos, mas não é validado automaticamente pela IA --- não existe verificação de que uma tarefa gerada testa efetivamente a competência DigComp~2.2 pretendida, e não apenas conhecimento genérico.

% ─────────────────────────────────────────────────────────────────────────────
\section{Comparação com o Estado da Arte} \label{sec:comparacao-arte}

O Play4Change insere-se no cruzamento de três linhas de trabalho revistas no Capítulo~2: plataformas gamificadas, sistemas de tutoria inteligente (ITS) e geração de conteúdo educativo por LLM.

\textbf{Face a plataformas gamificadas estabelecidas.} O Duolingo usa repetição espaçada e um currículo pré-construído por especialistas humanos; o Play4Change não usa SRS, mas distingue-se pela capacidade de gerar conteúdo a partir de qualquer fonte documental. Nem a Khan Academy nem o Coursera estão alinhados com o DigComp~2.2 ou permitem que instituições adicionem conteúdo próprio com geração automática de tarefas.

\textbf{Face a sistemas de tutoria inteligente clássicos.} Os ITS tradicionais (Cognitive Tutor, ALEKS) usam modelos de utilizador explícitos (BKT, PFA); o Play4Change usa uma heurística mais simples --- classificação de padrão de erro em três categorias --- trivial mas pedagogicamente interpretável. A vantagem dos ITS clássicos é a precisão diagnóstica; a do Play4Change é a generalidade, funcionando em qualquer domínio sem dados de calibração histórica.

\textbf{Face a soluções baseadas em LLM para educação.} Usar ChatGPT ou Claude diretamente para tutoria é flexível mas não estruturado pedagogicamente: sem percurso definido, sem gamificação, sem progressão garantida e sem isolamento do input do utilizador. O Play4Change estrutura a progressão (tarefa~$\to$~struggle~$\to$~explicação~$\to$~retentativa), gamifica, garante qualidade por schema e deduplicação, e isola o input em mensagens tipadas separadas.

\textbf{Contributos originais.} Três aspetos distinguem o Play4Change: a reutilização semântica de ramos adaptativos por similaridade coseno, que melhora continuamente a eficiência de custo sem intervenção humana; o modelo de escalamento estruturado (tarefa~$\to$~struggle~$\to$~explicação IA~$\to$~retentativa), com transições de estado explícitas; e a baralhagem determinística por utilizador, que previne partilha de respostas sem eliminar a reprodutibilidade dentro da mesma sessão.

% ─────────────────────────────────────────────────────────────────────────────
\section{Trabalho Futuro} \label{sec:trabalho-futuro}

\textbf{Funcionalidades sociais e de retenção.} A sugestão mais recorrente na avaliação de usabilidade foi a adição de funcionalidades sociais --- lista de amigos, tabelas de classificação, rivalidade positiva --- com impacto direto na retenção diária, já que a declaração de envolvimento mais baixa (3,92) foi precisamente sobre o regresso diário. Conquistas e medalhas por nível de aprendizagem podem ser implementadas sobre a infraestrutura de \texttt{micro\_competences} já existente, tal como notificações comportamentais personalizadas junto da janela habitual de utilização do utilizador.

\textbf{Reconhecimento formal via Open Badges e créditos ECTS.} O sistema de crachás atual é interno e não portável. Mapear estas conquistas para o \textit{standard} europeu \textit{Open Badges} permitiria a sua exportação e verificação externa, podendo suportar, no contexto da rede U!REKA, o reconhecimento formal de micro-créditos ECTS em unidades curriculares de reduzida carga horária.

A Tabela~\ref{tab:trabalho-futuro} resume cinco direções adicionais, de menor prioridade estratégica que as duas acima mas tecnicamente viáveis sobre a arquitetura atual sem alterações estruturais.

\begin{table}[h!]
\centering
\caption{Direções de trabalho futuro de prioridade secundária}
\label{tab:trabalho-futuro}
\resizebox{\textwidth}{!}{%
\begin{tabular}{lp{10.5cm}}
\hline
\textbf{Direção} & \textbf{Descrição} \\
\hline
Explicações adaptadas ao móvel & A avaliação qualitativa revelou que as explicações IA em prosa são longas demais para ecrãs pequenos; configurar o comprimento em função do \texttt{audienceLevel} e do canal, ou gerar uma versão resumida expansível. \\
Personalização geográfica & Um participante sugeriu usar a localização do dispositivo para personalizar questões sobre cidades inteligentes ao município do utilizador --- viável com a arquitetura atual, bastando injetar metadados de localização no contexto de geração. \\
Profundidade variável de struggle & Substituir \texttt{MAX\_STRUGGLE\_DEPTH~=~1} por um valor configurável por tópico, complementado por um modelo de utilizador simples baseado no histórico de \texttt{wrongAttemptCount} e na distribuição de padrões de erro. \\
Análise de desempenho das questões & O sistema já instrumenta \texttt{tasks\_submitted\_total} por resultado e tópico; um painel de qualidade (taxa de acerto, tempo médio) sinalizaria para revisão questões com taxa de acerto $<$20\% ou $>$95\%. \\
Validação automática de alinhamento DigComp~2.2 & Um segundo prompt poderia classificar a competência DigComp~2.2 testada por cada tarefa gerada e verificar a correspondência com o módulo configurado, sinalizando incongruências para revisão humana. \\
Suporte nativo multi-plataforma & O Kotlin Multiplatform possibilita compilação nativa para iOS com código partilhado, mas o projeto atual entrega apenas o alvo Android. \\
\hline
\end{tabular}%
}
\end{table}

\textbf{Duas direções estratégicas para o futuro.} A segmentação de perfis de afinidade tecnológica observada nos testes de usabilidade (§\ref{sec:usab-qualitativa}) sugere duas direções não necessariamente exclusivas: aprofundar a vertente universitária junto do perfil tecnologicamente confortável, através de Open Badges e créditos ECTS na rede U!REKA; ou, para o perfil de menor afinidade tecnológica --- que é precisamente quem mais beneficiaria do reforço de competências digitais ---, acrescentar mecanismos de acompanhamento, orientação ou mediação humana, em vez de depender exclusivamente da auto-motivação do utilizador. A escolha entre estas direções, ou a exploração de ambas em paralelo, fica em aberto para trabalho futuro.
